\scanpage{10}
\mspage{8}
Mais quelle est \inserted{(pour $q$ fixé)} la structure de $k\{T\}$-module de $H^{\bullet\,q}_{\mathrm{DRpd}}(\unsure{\mathscr X|X|\mathcal X}_{\ast},k)$ ? \deleted{\unsure{On calcule ce}}\inserted{On voit que} pour le \uline{degré total} (égal au degré extérieur. \uline{plus} $q$) on trouve $H^0(\mathscr X_\ast,k)\otimes_k k\{T\}$ tronqué en degré $\geq q$, donc
\[
H^{\bullet\,q}_{\mathrm{DRpd}}(\mathscr X_{\ast},k)\simeq\tau_q\big(H^0(\mathscr X_{\ast},k)\otimes_k k\{T\}\big)[q]
\]
\margin{(translation des degrés par $-q$)}

Si on réindexe les \unsure{bidegrés} par le couple (degré total, degré extérieur)
\[
{H'}^{n,q}_{\mathrm{DRpd}}(\mathscr X_{\ast},k)=H^{\overbrace{n-q}^{p},q}_{\mathrm{DRpd}}(\mathscr X_{\ast},k)
\]
(donc la condition de degré $p,q\geq0$ devient $n\geq q\geq 0$, l'op. diff est de bidegré $(0,1)$ donc c'est un hom (\uline{homogène}) de $S$-modules gradués), on trouve
\[
{H'}^{\bullet,q}_{\mathrm{DRpd}}(\mathscr X_\ast,k)\simeq\tau_q\big(H^q(\mathscr X_\ast,k)\otimes_k k\{T\}\big)
\]
Ces isom. sont compatibles avec les structures multiplicatives (et bien entendu fonctoriels en $\mathscr X_{\ast},k$ \ldots). Donc \inserted{\uline{a priori}} on ne récupère \inserted{(via $\mathord{H'}^{\bullet\bullet}_{\mathrm{DRpd}}(\mathscr X_{\ast},k)$)} pas les $k$-modules $H^q(\mathscr X_{\ast},k)$ et leurs \inserted{cup-}accouplements, mais les $k$-modules \inserted{gradués} $\tau_q(H^q(\mathscr X_{\ast},k)\otimes_k k\{T\})$ et leurs accouplements déduits du cup des $H^\bullet(\mathscr X_{\ast},k)$,
% archive sheet number: 9
